\section{Probability bounds for the cells}\label{sc:section}
Throughout this section, we work with the putative extremizer above:
\(C>B^2\), \(m\ge5\), and \(B=3/(2\sqrt{2\pi})\).
We will obtain lower bounds for each cell probability in terms of its
normalized centroid length. Since the probabilities sum to one, these
bounds constrain how the normalized centroid lengths can be
distributed among the cells.
Every terminating decimal in the proof denotes an exact rational number.
Write
\[
 \phi(x)=\frac{e^{-x^2/2}}{\sqrt{2\pi}},\qquad
 \Phi(x)=\int_{-\infty}^x\phi(u)\,du,\qquad
 r_i=\frac{|z_i|}{\sqrt C},\qquad P_i=\gamma_{m-1}(K_i).
\]
The winning cone \(K_i\) lies in \(\{x:\langle z_i,x\rangle\ge0\}\),
because the scores sum to zero. Thus \(0<P_i\le1/2\).

\begin{lemma}\label{sc:bounds}
Set
\[
 A=\begin{cases}0.929,&m=5,\\0.884,&m\ge6.\end{cases}
\]
For every \(i\),
\begin{equation}\label{sc:three-bounds}
 r_i\le\frac23,\qquad P_i\ge Ar_i^2,\qquad
 P_i\ge0.415r_i+0.15r_i^2.
\end{equation}
\end{lemma}

\subsection{Rearrangement bounds}
Fix a cell, and temporarily suppress its index.
Let \(e=z/|z|\), \(x=\Phi^{-1}(1-P)\), and
\(H=\{g:\langle e,g\rangle\ge x\}\). Since \(\gamma_{m-1}(H)=P\),
\[
 |z|-\int_H\langle e,g\rangle\,d\gamma_{m-1}(g)
 =\int(\langle e,g\rangle-x)(1_K-1_H)\,d\gamma_{m-1}(g)\le0.
\]
Consequently \(|z|\le\phi(x)\). As \(\sqrt C>B=(3/2)\phi(0)\),
this proves \(r<2/3\) and
\begin{equation}\label{sc:halfspace-ratio}
 \frac P{r^2}\ge \frac94 e^{x^2}\Phi(-x).
\end{equation}
For completeness, the minimum on the right can be localized by a
one-dimensional argument. Put \(M(x)=\phi(x)/\Phi(-x)\).
The logarithmic derivative of \(f(x)=(9/4)e^{x^2}\Phi(-x)\) is
\(2x-M(x)\), and
\[
 M'(x)=1-\operatorname{Var}(Z\mid Z>x)<1
\]
for a standard normal variable \(Z\). Hence \((\log f)'\) is strictly
increasing. Its signs at \(0.6119\) and \(0.6121\) are opposite, so
\(f\) has a unique global minimum in that interval.
The elementary bounds in Appendix~\ref{cert:arithmetic} give
\begin{equation}\label{sc:halfspace-minimum}
 \min_{x\in\mathbb R}f(x)>0.8843837>0.884.
\end{equation}
This proves the required quadratic estimate when \(m\ge6\).

When \(m=5\), conicality improves the estimate. Work in the
four-dimensional span of the centroids and write the standard Gaussian
as \(RU\), where \(U\) is uniform on \(S^3\) and independent of \(R\).
The coordinate \(\langle e,U\rangle\) has density
\((2/\pi)\sqrt{1-u^2}\) on \([-1,1]\). The same threshold comparison as
above shows that, among spherical sets of measure \(P\), an upper cap
maximizes the first moment in direction \(e\). Since \(P\le1/2\), write
its threshold as \(x\in[0,1)\). Its probability is
\[
 p(x)=\frac12-\frac{\arcsin x+x\sqrt{1-x^2}}\pi.
\]
Using \(\mathbb ER=3\sqrt{2\pi}/4\), its Gaussian first moment is
\[
 \mathbb ER\,\frac2\pi\int_x^1 u\sqrt{1-u^2}\,du
 =\phi(0)(1-x^2)^{3/2}.
\]
It follows that
\begin{equation}\label{sc:cap-ratio}
 \frac P{r^2}\ge R(x):=\frac94\frac{p(x)}{(1-x^2)^3}.
\end{equation}
To identify the global minimum, define
\[
 F(x)=6xp(x)-\frac2\pi(1-x^2)^{3/2},\qquad
 H(x)=p(x)-\frac x\pi\sqrt{1-x^2}.
\]
Direct differentiation gives
\[
 R'(x)=\frac94\frac{F(x)}{(1-x^2)^4},\qquad
 F'(x)=6H(x),\qquad
 H'(x)=\frac{4x^2-3}{\pi\sqrt{1-x^2}}.
\]
Thus \(H\) first decreases and then increases. Since \(H(0)=1/2\)
and \(H(1)=0\), it has exactly one zero in \((0,1)\); it is negative
after that zero. Accordingly \(F\) first increases and then decreases
to \(F(1)=0\). The certified signs \(F(0.29)<0<F(0.294)\) show that
\(F\) has exactly one interior zero and that this zero lies in
\((0.29,0.294)\). It is the unique minimizer of \(R\).
Appendix~\ref{cert:arithmetic} gives \(R(0.29)>0.9312526\) and
\(|R'|<0.03\) on this interval, whence
\begin{equation}\label{sc:cap-minimum}
 \min_{0\le x<1}R(x)>0.9311326>0.929.
\end{equation}
The first two assertions of Lemma~\ref{sc:bounds} follow.

\subsection{The loss on removing one score}
Rearrangement has supplied the radius bound and the quadratic
probability bound. To obtain the remaining linear--quadratic estimate
in Lemma~\ref{sc:bounds}, we use the behavior of the
winning probability under translation. Fix \(i\), put
\(X_j=\langle z_j,g\rangle\), and set \(\alpha=m/(m-1)\).
Removing \(z_i\) and recentering the remaining vectors produces
\(z_j+z_i/(m-1)\), \(j\ne i\), whose squared norms sum to
\(C-\alpha|z_i|^2\). A common translation of the scores has zero
expectation. The extremal dual bound therefore gives
\[
 \mathbb E\max_{j\ne i}X_j\le
 \sqrt C\sqrt{C-\alpha|z_i|^2}.
\]
Since \(\mathbb E\max_jX_j=C\), the loss satisfies
\begin{equation}\label{sc:loss-lower}
 \mathbb E\bigl(X_i-\max_{j\ne i}X_j\bigr)_+
 \ge C\bigl(1-\sqrt{1-\alpha r_i^2}\bigr).
\end{equation}
Here \(u_+=\max(u,0)\).

The vectors \(z_i-z_j\), \(j\ne i\), form a basis of the centroid
span. There is therefore a vector \(h\) in that span satisfying
\(\langle h,z_i-z_j\rangle=1\) for every \(j\ne i\).
For \(u\ge0\), the gap probability is
\[
 p_i(u):=\Pr\{X_i\ge X_j+u\text{ for all }j\ne i\}
       =\gamma_{m-1}(K_i+uh).
\]
Summing the equations defining \(h\) gives
\(\langle h,z_i\rangle=(m-1)/m=1/\alpha\).
Differentiation of the translated Gaussian density, justified by
Gaussian integrability, yields
\begin{equation}\label{sc:translation-derivative}
 p_i'(0)=-\int_{K_i}\langle h,x\rangle\,d\gamma_{m-1}(x)
        =-\langle h,z_i\rangle=-\frac1\alpha.
\end{equation}

Ehrhard's inequality states that \(\Phi^{-1}\circ\gamma_{m-1}\) is
concave under Minkowski interpolation of convex sets
\cite{Ehrhard}; we use the formulation in
\cite[Theorem~1.1]{Borell}, which also allows general Borel sets.
The cone \(K_i\) is closed and convex. Its score-difference normals
form a basis, so it has nonempty interior and is a proper subset of
the centroid span. Minkowski interpolation of two translates of \(K_i\) is the translate
with the interpolated displacement. Hence
\(u\mapsto\Phi^{-1}(p_i(u))\) is concave. All finite translates
have measure strictly between zero and one, and the derivative at
zero exists by \eqref{sc:translation-derivative}. With
\(q=\Phi^{-1}(P_i)\), the supporting tangent gives
\[
 p_i(u)\le\Phi\!\left(q-\frac{u}{\alpha\phi(q)}\right)
 \qquad(u\ge0).
\]
Integration, using \((q\Phi(q)+\phi(q))'=\Phi(q)\), gives
\begin{equation}\label{sc:loss-upper}
 \mathbb E\bigl(X_i-\max_{j\ne i}X_j\bigr)_+
 =\int_0^\infty p_i(u)\,du
 \le\alpha\phi(q)\bigl(q\Phi(q)+\phi(q)\bigr).
\end{equation}
Define
\[
 G(q)=\frac{\phi(q)\bigl(q\Phi(q)+\phi(q)\bigr)}{\Phi(q)^2}.
\]
Combining \eqref{sc:loss-lower} and \eqref{sc:loss-upper}, rationalizing,
and canceling \(\alpha\), we obtain
\begin{equation}\label{sc:loss-constraint}
 P_i^2G\bigl(\Phi^{-1}(P_i)\bigr)
 \ge\frac{Cr_i^2}{1+\sqrt{1-\alpha r_i^2}}
 \ge\frac{B^2r_i^2}{1+\sqrt{1-r_i^2}}.
\end{equation}

\subsection{A linear--quadratic probability bound}
The loss constraint is implicit in the cell probability. To extract
an explicit lower bound, we first prove two monotonicity properties
of \(G\).
\begin{lemma}\label{sc:truncated-variance}
The function \(G\) is nonincreasing and satisfies \(0<G(q)<1\).
Moreover,
\[
 P\longmapsto P^2G\bigl(\Phi^{-1}(P)\bigr)
\]
is increasing on \((0,1/2]\).
\end{lemma}
\begin{proof}
Put \(\lambda=\phi(q)/\Phi(q)\). Integration of the normal density gives
\[
 \mathbb E(Z\mid Z\le q)=-\lambda,\qquad
 V(q):=\operatorname{Var}(Z\mid Z\le q)=1-q\lambda-\lambda^2.
\]
Thus \(G=1-V=\lambda(q+\lambda)\). The distance
\(U=q-Z\), under this conditioning, has positive mean
\(q+\lambda\) and positive variance, proving \(0<G<1\).
Its survival function is
\(S(u)=\Phi(q-u)/\Phi(q)\), \(u\ge0\).
The identity
\((\log\Phi)''(q)=-\lambda(q+\lambda)<0\)
shows that \(\log S\) is concave. Its increments therefore satisfy
\[
 \log S(u+v)-\log S(u)\le\log S(v)-\log S(0).
\]
Since \(S(0)=1\), this gives \(S(u+v)\le S(u)S(v)\).
Integrating over \(u,v\ge0\) gives
\[
 \frac12\mathbb EU^2
 =\int_0^\infty\!\int_0^\infty S(u+v)\,du\,dv
 \le(\mathbb EU)^2.
\]
Hence \(V\le(q+\lambda)^2\). Direct differentiation now gives
\[
 V'(q)=\lambda\bigl((q+\lambda)^2-V(q)\bigr)\ge0,
\]
so \(G\) is nonincreasing.
Finally, for \(q\le0\), both positive factors of
\(\phi(q)(q\Phi(q)+\phi(q))\) are increasing in \(q\), proving the
last assertion.
\end{proof}

We finish the proof of Lemma~\ref{sc:bounds}. Set
\(P_0(r)=(0.415+0.15r)r\) for \(0<r\le2/3\). Then
\(P_0(r)\le P_0(2/3)=103/300<1/2\).
We will prove
\begin{equation}\label{sc:test-probability}
 (0.415+0.15r)^2G\bigl(\Phi^{-1}(P_0(r))\bigr)
 <\frac{B^2}{1+\sqrt{1-r^2}}.
\end{equation}
By Lemma~\ref{sc:truncated-variance} and
\eqref{sc:loss-constraint}, this implies \(P_i>P_0(r_i)\), which is
stronger than the last assertion of Lemma~\ref{sc:bounds}.

Here is a finite interval verification of \eqref{sc:test-probability}.
In each row of Table~\ref{sc:interval-table}, \(U\) bounds
\(G(\Phi^{-1}(P_0(r)))\) above throughout \([a,b]\cap(0,2/3]\).
The first row uses \(G<1\); the other rows follow from monotonicity and
the certified quantile bounds in Appendix~\ref{cert:arithmetic}.
The last column is a strict lower bound for
\[
 \Delta(a,b,U):=
 \frac{B^2}{1+\sqrt{1-a^2}}-U(0.415+0.15b)^2.
\]
It is positive in every row. The coefficient \((0.415+0.15r)^2\)
and the right side of \eqref{sc:test-probability} are increasing in
\(r\), so these endpoint inequalities cover every point of the
intervals, not only their endpoints.

\begin{table}[ht]
\centering
\begin{tabular}{cccc}
\hline
\(a\)&\(b\)&\(U\)&\(10^6\Delta(a,b,U)>\)\\
\hline
0&0.05&1&542\\
0.05&0.14&0.91&6173\\
0.14&0.24&0.87&2976\\
0.24&0.33&0.836&1328\\
0.33&0.41&0.809&523\\
0.41&0.48&0.788&392\\
0.48&0.54&0.771&1076\\
0.54&0.59&0.756&2787\\
0.59&0.64&0.743&4115\\
0.64&\(2/3\)&0.73&8886\\
\hline
\end{tabular}
\caption{Bounds proving \eqref{sc:test-probability}. All displayed
finite decimals are exact rational numbers.}
\label{sc:interval-table}
\end{table}
